% GENERATED FILE. Edit canonical lessons, then run npm run generate:books.
% canonical-source-hash: fnv1a64:315132eab7e40c78
% canonical-lessons: ZH-C42-dongxi, ZH-C42-mai-sell, ZH-C42-shangdian, ZH-C42-pianyi, ZH-C42-gui, ZH-R42-review-37-42

\chapter{In the Shop}
\label{ch:zh-in-the-shop}
\hlchaptermodality{42}

\begin{chapteropening}
\textbf{By the end of this chapter:} \emph{I can say what a shop sells and whether it is cheap or dear.}

Complete the last lesson of chapter 42: I can say what a shop sells and whether it is cheap or dear.
\end{chapteropening}

\section[dōngxi]{dōngxi — thing; things}

\label{lesson:ZH-C42-dongxi}



\begin{quote}
Before the new one: say the Mandarin for too, then the Mandarin for to want.
\end{quote}

\subsection*{You'll want to know: dōngxi}
\emph{dōngxi} — \textquotedblleft{}thing; things\textquotedblright{}.

\begin{quote}
\emph{nǐ mǎi shénme dōngxi} — What are you buying?
\end{quote}

Its characters come later; for now it is a word for the ear.

\begin{cousinweb}
\emph{dōngxi} is built from dōng 'east' and xī 'west', both inherited from Middle Chinese, \emph{tuwng} and \emph{sej} (Guangyun, in Baxter's transcription). How east-and-west came to mean a thing is not settled.
\end{cousinweb}

\begin{grammarlens}[title={What you've built}]
One more word for asking politely.
\end{grammarlens}

\subsection*{Your turn}
\begin{itemize}
\raggedright
  \item \emph{Say it:} \emph{dōngxi}
  \item \emph{Say it:} \emph{dōngxi}, once more
  \item \emph{Say it:} \emph{nǐ mǎi shénme dōngxi} — What are you buying?
  \item \emph{Recall:} say the Mandarin for too, then the Mandarin for to want, then say \emph{dōngxi} again
\end{itemize}

\begin{tcolorbox}[breakable,colback=teal!4,colframe=teal!35!black,title={Before you move on}]
What does dōngxi mean? (\textquotedblleft{}Thing; things\textquotedblright{}.) Say it once more.
\end{tcolorbox}
\section[mài]{mài — to sell}

\label{lesson:ZH-C42-mai-sell}



\begin{quote}
Before the new one: say the Mandarin for to want, then the Mandarin for thing.
\end{quote}

\subsection*{You'll want to know: mài}
\emph{mài} — \textquotedblleft{}to sell\textquotedblright{}.

\begin{quote}
\emph{tāmen mài shénme} — What do they sell?
\end{quote}

Its characters come later; for now it is a word for the ear.

\begin{cousinweb}
\emph{mài} is inherited from Middle Chinese \emph{meaH} (Guangyun, in Baxter's transcription), glossed 'to put goods out'. It differs from the word for buy only in its tone.
\end{cousinweb}

\begin{grammarlens}[title={What you've built}]
One more word for asking politely.
\end{grammarlens}

\subsection*{Your turn}
\begin{itemize}
\raggedright
  \item \emph{Say it:} \emph{mài}
  \item \emph{Say it:} \emph{mài}, once more
  \item \emph{Say it:} \emph{tāmen mài shénme} — What do they sell?
  \item \emph{Recall:} say the Mandarin for to want, then the Mandarin for thing, then say \emph{mài} again
\end{itemize}

\begin{tcolorbox}[breakable,colback=teal!4,colframe=teal!35!black,title={Before you move on}]
What does mài mean? (\textquotedblleft{}To sell\textquotedblright{}.) Say it once more.
\end{tcolorbox}
\section[shāngdiàn]{shāngdiàn — shop}

\label{lesson:ZH-C42-shangdian}



\begin{quote}
Before the new one: say the Mandarin for thing, then the Mandarin for to sell.
\end{quote}

\subsection*{You'll want to know: shāngdiàn}
\emph{shāngdiàn} — \textquotedblleft{}shop\textquotedblright{}.

\begin{quote}
\emph{zhège shāngdiàn hěn xiǎo} — This shop is tiny.
\end{quote}

Its characters come later; for now it is a word for the ear.

\begin{cousinweb}
\emph{shāngdiàn} is built from shāng, usually explained as 'trade', and diàn 'shop', both inherited from Middle Chinese, \emph{syang} and \emph{temH} (Guangyun, in Baxter's transcription); the dictionary explains the second as a place where goods are kept and sold.
\end{cousinweb}

\begin{grammarlens}[title={What you've built}]
One more word for asking politely.
\end{grammarlens}

\subsection*{Your turn}
\begin{itemize}
\raggedright
  \item \emph{Say it:} \emph{shāngdiàn}
  \item \emph{Say it:} \emph{shāngdiàn}, once more
  \item \emph{Say it:} \emph{zhège shāngdiàn hěn xiǎo} — This shop is tiny.
  \item \emph{Recall:} say the Mandarin for thing, then the Mandarin for to sell, then say \emph{shāngdiàn} again
\end{itemize}

\begin{tcolorbox}[breakable,colback=teal!4,colframe=teal!35!black,title={Before you move on}]
What does shāngdiàn mean? (\textquotedblleft{}Shop\textquotedblright{}.) Say it once more.
\end{tcolorbox}
\section[piányi]{piányi — cheap}

\label{lesson:ZH-C42-pianyi}



\begin{quote}
Before the new one: say the Mandarin for to sell, then the Mandarin for shop.
\end{quote}

\subsection*{You'll want to know: piányi}
\emph{piányi} — \textquotedblleft{}cheap\textquotedblright{}.

\begin{quote}
\emph{zhège hěn piányi} — This one's cheap.
\end{quote}

Its characters come later; for now it is a word for the ear.

\begin{cousinweb}
\emph{piányi} is built from two old words, both inherited from Middle Chinese: \emph{bjien} 'at ease' and \emph{ngje} 'fitting' (Guangyun, in Baxter's transcription). The sense 'cheap' is not given there, so how the pair came to mean it is not settled.
\end{cousinweb}

\begin{grammarlens}[title={What you've built}]
One more word for asking politely.
\end{grammarlens}

\subsection*{Your turn}
\begin{itemize}
\raggedright
  \item \emph{Say it:} \emph{piányi}
  \item \emph{Say it:} \emph{piányi}, once more
  \item \emph{Say it:} \emph{zhège hěn piányi} — This one's cheap.
  \item \emph{Recall:} say the Mandarin for to sell, then the Mandarin for shop, then say \emph{piányi} again
\end{itemize}

\begin{tcolorbox}[breakable,colback=teal!4,colframe=teal!35!black,title={Before you move on}]
What does piányi mean? (\textquotedblleft{}Cheap\textquotedblright{}.) Say it once more.
\end{tcolorbox}
\section[guì]{guì — expensive}

\label{lesson:ZH-C42-gui}



\begin{quote}
Before the new one: say the Mandarin for shop, then the Mandarin for cheap.
\end{quote}

\subsection*{You'll want to know: guì}
\emph{guì} — \textquotedblleft{}expensive\textquotedblright{}.

\begin{quote}
\emph{nàge tài guì le} — That one's too expensive.
\end{quote}

Its characters come later; for now it is a word for the ear.

\begin{cousinweb}
\emph{guì} is inherited from Middle Chinese \emph{kjw+jH} (Guangyun, in Baxter's transcription), glossed 'honoured, high'. The price sense is not given there; from valued to costly is a likely step, but it is not settled.
\end{cousinweb}

\begin{grammarlens}[title={What you've built}]
One more word for asking politely.
\end{grammarlens}

\subsection*{Your turn}
\begin{itemize}
\raggedright
  \item \emph{Say it:} \emph{guì}
  \item \emph{Say it:} \emph{guì}, once more
  \item \emph{Say it:} \emph{nàge tài guì le} — That one's too expensive.
  \item \emph{Recall:} say the Mandarin for shop, then the Mandarin for cheap, then say \emph{guì} again
\end{itemize}

\begin{tcolorbox}[breakable,colback=teal!4,colframe=teal!35!black,title={Before you move on}]
What does guì mean? (\textquotedblleft{}Expensive\textquotedblright{}.) Say it once more.
\end{tcolorbox}
\section[(dialogue)]{Looking back: from This One and That One to In the Shop}

\label{lesson:ZH-R42-review-37-42}



\begin{quote}
Say the Mandarin for this and this one.
\end{quote}

\subsection*{Your turn}
\begin{itemize}
\raggedright
  \item \emph{Say it:} this, this one, that, which one, that one
  \item \emph{Say it:} to eat, meal, cooked rice, hungry, dish of food
  \item \emph{Say it:} to drink, water, thirsty, milk, tea
  \item \emph{Say it:} to buy, apple, fruit, bread, egg
  \item \emph{Say it:} how much, money, yuan, too, to want
  \item \emph{Say it:} thing, to sell, shop, cheap, expensive
  \item \emph{Read it:} \zh{这} \zh{那} \zh{吃} \zh{饭} \zh{喝} \zh{水} \zh{买} \zh{多} \zh{少} — each one aloud, then a word that holds it
\end{itemize}

\begin{tcolorbox}[breakable,colback=teal!4,colframe=teal!35!black,title={Before you move on}]
This? (\textbf{\zh{这}}, \emph{zhè}.) This one? (\textbf{\zh{这个}}, \emph{zhège}.) That? (\textbf{\zh{那}}, \emph{nà}.) Which one? (\emph{nǎge}.) That one? (\textbf{\zh{那个}}, \emph{nàge}.) To eat? (\textbf{\zh{吃}}, \emph{chī}.) Meal? (\textbf{\zh{饭}}, \emph{fàn}.) Cooked rice? (\emph{mǐfàn}.) Hungry? (\emph{è}.) Dish of food? (\emph{cài}.) To drink? (\textbf{\zh{喝}}, \emph{hē}.) Water? (\textbf{\zh{水}}, \emph{shuǐ}.) Thirsty? (\emph{kě}.) Milk? (\emph{niúnǎi}.) Tea? (\emph{chá}.) To buy? (\textbf{\zh{买}}, \emph{mǎi}.) Apple? (\emph{píngguǒ}.) Fruit? (\emph{shuǐguǒ}.) Bread? (\emph{miànbāo}.) Egg? (\emph{jīdàn}.) How much? (\textbf{\zh{多少}}, \emph{duōshao}.) Money? (\emph{qián}.) Yuan? (\emph{kuài}.) Too? (\emph{tài}.) To want? (\emph{yào}.) Thing? (\emph{dōngxi}.) To sell? (\emph{mài}.) Shop? (\emph{shāngdiàn}.) Cheap? (\emph{piányi}.) Expensive? (\emph{guì}.)
\end{tcolorbox}
